% © 2025 Ing. David Jaroš — CC BY-NC-ND 4.0
%
% This work is licensed under a Creative Commons Attribution-NonCommercial-NoDerivatives 
% 4.0 International License (CC BY-NC-ND 4.0).
%
% License History: Earlier drafts (up to v0.3) were released under CC BY 4.0. 
% From v0.4 onward, all material is released under CC BY-NC-ND 4.0 to protect 
% the integrity of the theoretical work during ongoing academic development.
%
% See LICENSE.md for full license text.

% This file can be compiled standalone or included in another document
\ifdefined\INCLUDEMODE
  % Being included in another document - skip preamble
\else
  % Standalone compilation - include preamble
  \documentclass[12pt]{article}
  \usepackage[a4paper, margin=2.5cm]{geometry}
  \usepackage{amsmath, amssymb, amsthm}
  \usepackage{hyperref}
  \usepackage{xcolor}
  \usepackage{mdframed}
  \usepackage{slashed}

  \newtheorem{theorem}{Theorem}[section]
  \newtheorem{lemma}[theorem]{Lemma}
  \newtheorem{proposition}[theorem]{Proposition}
  \theoremstyle{definition}
  \newtheorem{definition}[theorem]{Definition}
  \theoremstyle{remark}
  \newtheorem{remark}[theorem]{Remark}

  \newmdenv[backgroundcolor=yellow!20, linecolor=orange, linewidth=1.5pt]{gapbox}
  \newmdenv[backgroundcolor=green!15, linecolor=green!60!black, linewidth=1.5pt]{resultbox}

  \title{\textbf{Step 3: Dirac Structure from the Biquaternion Algebra $\mathbb{C}\otimes\mathbb{H}$}}
  \author{David Jaroš}
  \date{2026-03-06}

  \begin{document}
  \maketitle

  \begin{abstract}
  We verify the claim that the Dirac gamma matrices and the Dirac equation emerge from
  the biquaternion algebra $\mathbb{B} = \mathbb{C}\otimes\mathbb{H} \cong \mathrm{Mat}(2,\mathbb{C})$.
  \textbf{Verdict: SKETCH.}
  The algebra isomorphism $\mathbb{C}\otimes\mathbb{H} \cong \mathrm{Mat}(2,\mathbb{C})$ and the
  identification of Pauli matrices with quaternion generators is a \emph{proved} algebraic fact.
  The construction of Dirac gamma matrices via the tensor product $\mathbb{B}\otimes\mathbb{B}$
  is structurally complete.  However, the derivation of the massive Dirac equation from the
  Euler--Lagrange equations of $S[\Theta]$ contains a gap: the mass term requires
  constraining $\Theta$ to the spinorial subspace plus an identification of the mass parameter,
  neither of which is derived from $S[\Theta]$ without additional assumptions.
  \end{abstract}
\fi

\section{Step 3: Dirac Structure from $\mathbb{C}\otimes\mathbb{H}$}
\label{sec:dirac_emergence}

\subsection{The Algebra Isomorphism}

\begin{theorem}[Standard Algebra Isomorphism]
\label{thm:biquat_iso}
As algebras over $\mathbb{R}$:
\[
  \mathbb{B} \;=\; \mathbb{C}\otimes_\mathbb{R}\mathbb{H} \;\cong\; \mathrm{Mat}(2,\mathbb{C}).
\]
\end{theorem}

\begin{proof}[Proof (standard; included for completeness)]
The quaternion basis $\{1, \mathbf{i}, \mathbf{j}, \mathbf{k}\}$ admits the $2\times 2$
complex matrix representation:
\begin{equation}
  1 \mapsto \begin{pmatrix}1&0\\0&1\end{pmatrix},\quad
  \mathbf{i} \mapsto \begin{pmatrix}i&0\\0&-i\end{pmatrix},\quad
  \mathbf{j} \mapsto \begin{pmatrix}0&1\\-1&0\end{pmatrix},\quad
  \mathbf{k} \mapsto \begin{pmatrix}0&i\\i&0\end{pmatrix}.
  \label{eq:quat_to_mat}
\end{equation}
These satisfy $\mathbf{i}^2=\mathbf{j}^2=\mathbf{k}^2=\mathbf{ijk}=-1$ and generate
$\mathrm{Mat}(2,\mathbb{C})$ (dimension $4$ over $\mathbb{C}$, or $8$ over $\mathbb{R}$).
Since $\mathbb{B} = \mathbb{C}\otimes\mathbb{H}$ also has dimension $8$ over $\mathbb{R}$
and is a simple algebra, the map is an isomorphism.
\end{proof}

\begin{resultbox}
\textbf{Established (proved)}: $\mathbb{C}\otimes\mathbb{H}\cong\mathrm{Mat}(2,\mathbb{C})$,
with generators mapping to $i\sigma_1, i\sigma_2, i\sigma_3$ (up to signs), where $\sigma_a$
are the Pauli matrices.
\end{resultbox}

\subsection{Pauli Matrices from Biquaternion Generators}

Under the map \eqref{eq:quat_to_mat}, the imaginary-unit quaternion generators satisfy:
\begin{align}
  \mathbf{i}\mapsto i\sigma_3,\quad
  \mathbf{j}\mapsto i\sigma_2,\quad
  \mathbf{k}\mapsto i\sigma_1,
\end{align}
where
\begin{equation}
  \sigma_1 = \begin{pmatrix}0&1\\1&0\end{pmatrix},\quad
  \sigma_2 = \begin{pmatrix}0&-i\\i&0\end{pmatrix},\quad
  \sigma_3 = \begin{pmatrix}1&0\\0&-1\end{pmatrix}.
\end{equation}

\begin{resultbox}
\textbf{Established (proved)}: Pauli matrices $\sigma_1,\sigma_2,\sigma_3$ are the generators
of the biquaternion algebra $\mathbb{B}$ under the isomorphism of Theorem~\ref{thm:biquat_iso}.
They are \emph{not postulated}; they arise from the quaternion multiplication rules.
\end{resultbox}

\subsection{Dirac Matrices from $\mathbb{B}\otimes\mathbb{B}$}

The Dirac algebra requires a $4\times 4$ representation (Clifford algebra $\mathrm{Cl}_{1,3}$).
This arises from the tensor product:
\begin{equation}
  \mathrm{Mat}(2,\mathbb{C})\otimes\mathrm{Mat}(2,\mathbb{C}) \;\cong\; \mathrm{Mat}(4,\mathbb{C}).
  \label{eq:tensor_product}
\end{equation}

The standard Weyl representation of the Dirac matrices is:
\begin{equation}
  \gamma^0 = \begin{pmatrix}0 & \mathbb{1}_2\\\mathbb{1}_2 & 0\end{pmatrix},\qquad
  \gamma^i = \begin{pmatrix}0 & \sigma_i\\-\sigma_i & 0\end{pmatrix},\quad i=1,2,3.
  \label{eq:weyl_gamma}
\end{equation}

These are tensor products of Pauli matrices:
\begin{equation}
  \gamma^0 = \sigma_1\otimes\mathbb{1}_2,\qquad
  \gamma^i = i\sigma_2\otimes\sigma_i,\quad i=1,2,3.
  \label{eq:gamma_tensor}
\end{equation}

\begin{proposition}[Gamma Matrices from Biquaternion Algebra]
The Dirac gamma matrices in the Weyl representation are tensor products of the generators
of $\mathbb{B}\cong\mathrm{Mat}(2,\mathbb{C})$.  Explicitly, via \eqref{eq:gamma_tensor},
each $\gamma^\mu$ lies in $\mathbb{B}\otimes\mathbb{B}\cong\mathrm{Mat}(4,\mathbb{C})$.
\end{proposition}

\begin{proof}
Direct computation using \eqref{eq:tensor_product}.  The Clifford relation
$\{\gamma^\mu,\gamma^\nu\}=2g^{\mu\nu}\mathbb{1}_4$ is verified by the identity
$\{\sigma_i,\sigma_j\}=2\delta_{ij}$ and the anticommutation of Pauli matrices.
\end{proof}

\begin{resultbox}
\textbf{Established (proved)}: The Dirac gamma matrices in Weyl representation arise from
the tensor product of two copies of the biquaternion algebra.  No additional structure is
needed beyond $\mathbb{C}\otimes\mathbb{H}$.
\end{resultbox}

\subsection{The Dirac-like Operator in UBT}

From \texttt{ubt/operators/dirac\_like\_operator.tex} (Theorem~2.1 in that document):

\begin{theorem}[UBT Dirac Operator; from \texttt{dirac\_like\_operator.tex}]
The operator
\[
  \mathcal{D}\Theta = i\gamma^\mu\nabla_\mu\Theta
\]
acting on the biquaternionic bundle $\mathcal{E} = P\times_{G_{\rm tot}}(\mathbb{B}\otimes S\otimes V_G)$
is uniquely determined (up to gauge/unitary equivalence) by requiring:
\begin{enumerate}
  \item Clifford algebra structure: symbol $\sigma(\mathcal{D})=\gamma^\mu\xi_\mu$,
  \item Gauge covariance: $[\mathcal{D}, U]=0$ for $U\in G_{\rm gauge}$,
  \item Metric compatibility.
\end{enumerate}
\end{theorem}

The proof in \texttt{dirac\_like\_operator.tex} is complete and uses only properties of
the Levi-Civita theorem plus gauge theory.

\subsection{Does the Dirac Equation Follow from $S[\Theta]$?}

The Dirac equation $(i\gamma^\mu\partial_\mu - m)\psi = 0$ follows from varying the
Dirac action:
\begin{equation}
  S_{\rm Dirac} = \int\bar\psi(i\gamma^\mu\partial_\mu - m)\psi\,d^4x.
\end{equation}

\begin{gapbox}
\textbf{Gap D1 — Spinorial Subspace Constraint Not Derived}:
In the UBT action $S[\Theta]$, the field $\Theta$ takes values in the full biquaternionic
bundle $\mathbb{B}\otimes S\otimes V_G$.  The Dirac equation acts on the spinor subspace $S$.

Appendix~FORMAL (Proposition 2) states: ``when $\Theta$ is constrained to the spinorial
subspace\ldots the linearized field equation becomes the Dirac equation.''

This constraint is \emph{not derived} from $S[\Theta]$.  It must be imposed by hand.
Without it, the Euler-Lagrange equations of $S[\Theta]$ give a more general equation.

\textbf{Status}: OPEN sub-task.
\end{gapbox}

\begin{gapbox}
\textbf{Gap D2 — Mass Term Origin}:
The mass term $m\bar\psi\psi$ in the Dirac action requires an explicit source.
In UBT, the mass of the electron/fermion is a separate derivation problem
(addressed in \texttt{appendix\_E\_m0\_derivation\_strict.tex}).

The Dirac \emph{equation} (massless) $i\gamma^\mu\partial_\mu\psi=0$ follows naturally
from $\mathcal{D}\Theta = 0$ in the spinorial sector.  The massive case requires identifying
a UBT-derived mass parameter.

\textbf{Status}: OPEN sub-task (linked to fermion mass problem).
\end{gapbox}

\subsection{Summary of What Is Established}

\begin{center}
\renewcommand{\arraystretch}{1.3}
\begin{tabular}{lll}
\hline
\textbf{Result} & \textbf{Status} & \textbf{Source} \\
\hline
$\mathbb{C}\otimes\mathbb{H}\cong\mathrm{Mat}(2,\mathbb{C})$ (algebra iso.)
  & \textbf{Proved} & Standard algebra \\
Pauli matrices $\sigma_i$ from quaternion generators
  & \textbf{Proved} & Theorem~\ref{thm:biquat_iso} \\
$\gamma^\mu$ from $\mathbb{B}\otimes\mathbb{B}$ tensor products
  & \textbf{Proved} & \eqref{eq:gamma_tensor} \\
$\mathcal{D} = i\gamma^\mu\nabla_\mu$ uniquely determined
  & \textbf{Proved} & \texttt{dirac\_like\_operator.tex} Thm.~2.1 \\
Massless Dirac eq.\ from $\mathcal{D}\Theta=0$ in spinorial sector
  & \textbf{Sketch} & Gap D1 (constraint not derived) \\
Massive Dirac eq.\ from $S[\Theta]$
  & \textbf{Sketch} & Gap D2 (mass term) \\
\hline
\end{tabular}
\end{center}

\subsection{Verdict}

\begin{center}
\large\textbf{VERDICT: SKETCH}
\end{center}

The \emph{algebraic} emergence of Dirac structure from $\mathbb{C}\otimes\mathbb{H}$ is
\textbf{fully proved}:
\begin{itemize}
  \item $\mathbb{C}\otimes\mathbb{H}\cong\mathrm{Mat}(2,\mathbb{C})$ (standard algebra),
  \item Pauli matrices arise from quaternion generators,
  \item Dirac $\gamma^\mu$ arise from $\mathbb{B}\otimes\mathbb{B}$,
  \item The Dirac-like operator $\mathcal{D}$ is uniquely determined.
\end{itemize}

The \emph{dynamical} emergence of the Dirac \emph{equation} from $S[\Theta]$ is a
\textbf{Sketch}, blocked by:
\begin{enumerate}
  \item Deriving the spinorial subspace constraint from $S[\Theta]$ (Gap D1),
  \item Providing a UBT-derived mass term (Gap D2).
\end{enumerate}

\noindent\textbf{Source files checked}:
\begin{itemize}
  \item \texttt{ubt/operators/dirac\_like\_operator.tex}
  \item \texttt{consolidation\_project/appendix\_FORMAL\_qm\_gr\_unification.tex}
\end{itemize}
\noindent\textbf{Date}: 2026-03-06\\
\textbf{Author}: David Jaroš

\ifdefined\INCLUDEMODE
\else
\end{document}
\fi
